\emph{Step 1 (the number Gram).} With $\Theta_{k}[(i,A),(j,B)] = \tfrac12\varepsilon_{kij}\varepsilon_{AB}$ as in Theorem \claimref{families.skeleton-rank-one-plus-ordering-shift}, $\Theta_{k}\Theta_{l} = \tfrac14 (E_{k}E_{l}) \otimes \varepsilon^{2} = -\tfrac14 E_{k}E_{l} \otimes \mathbb{1}$ with $(E_{k})_{ij} = \varepsilon_{kij}$, and $(E_{k}E_{l})_{im} = \sum_{j}\varepsilon_{kij}\varepsilon_{ljm} = \delta_{il}\delta_{km} - \delta_{kl}\delta_{im}$, so $\Theta_{k}\Theta_{l} = \tfrac14(\delta_{kl}\mathbb{1} - e_{l}e_{k}^{\mathsf{T}} \otimes \mathbb{1})$. Then $2\mathcal{N}_{kl} = 2\langle a^{\dagger}\Theta_{k}\Theta_{l}a\rangle = \tfrac12(\delta_{kl}n - \sum_{A}\langle a_{lA}^{\dagger}a_{kA}\rangle) = \tfrac12(\delta_{kl}n - \rho_{lk})$. The check verifies the six-by-six identity for all $k, l$ and the state-level identity on three exact non-Gaussian Fock superpositions. Since $\bar\omega\omega^{\mathsf{T}} \ge 0$ and $K \ge 0$ (Proposition \claimref{synthesis.face-mass-form-unit-plus-gram}(iii)), $H \ge \mathbb{1} + \tfrac12(n\mathbb{1} - \rho^{\mathsf{T}})$, whose least eigenvalue is $1 + \tfrac12(n - \rho_{\max})$ ($\rho^{\mathsf{T}}$ has the eigenvalues of $\rho$). In the skeleton, by Proposition \claimref{synthesis.face-mass-form-unit-plus-gram}(iii), $2\mathcal{N} = x(\mathbb{1} + P)$ with $P = P_{[\bar v]}$; comparing traces, $\tfrac12(3n - n) = 4x$, so $n = 4x$ and $\rho^{\mathsf{T}} = n\mathbb{1} - 2x(\mathbb{1} + P) = 2x(\mathbb{1} - P)$, with eigenvalues $2x, 2x, 0$: the heavy direction is empty and the bound is $1 + \tfrac12(4x - 2x) = 1 + x = \cosh^{2}r$, the light eigenvalue of the theorem. This is (i).

\emph{Step 2 (fixed second moments).} $\omega_{k} = \sum_{pq}\Theta_{k}[p,q]\langle a_{p}a_{q}\rangle$ and $\mathcal{N}_{kl} = \sum_{pq}(\Theta_{k}\Theta_{l})[p,q]\langle a_{p}^{\dagger}a_{q}\rangle$ depend only on the second moments, so for $\psi$ with the skeleton's one- and two-point functions they take the skeleton's values. In the skeleton, by the theorem, $\bar\omega\omega^{\mathsf{T}} + K + 2\mathcal{N} = H - \mathbb{1} = x\mathbb{1} + x(5 + 6x)P$, and by Proposition \claimref{synthesis.face-mass-form-unit-plus-gram}(iii) $K = 2x^{2}P$ and $2\mathcal{N} = x(\mathbb{1} + P)$; so $\bar\omega\omega^{\mathsf{T}} = (x(5 + 6x) - 2x^{2} - x)P = 4x(1 + x)P$ (the check verifies the polynomial identity and the trace identity $\tr H = 3 + n + \sum_{k}\|\hat\Omega_{k}\psi\|^{2}$ at the skeleton). Hence for $\psi$: $H(\psi) = \mathbb{1} + 4x(1 + x)P + K(\psi) + x(\mathbb{1} + P) = (1 + x)\mathbb{1} + (5x + 4x^{2})P + K(\psi)$, with $K(\psi) \ge 0$ because it is the Gram matrix of the vectors $(\hat\Omega_{k} - \omega_{k})\psi$. The last two terms are positive semidefinite, so $H(\psi) - (1 + x)\mathbb{1} \ge 0$; the check confirms the principal-minor criterion on sample Gram matrices. Connected cumulants of every order (and products of the mean with third cumulants, when the state is displaced) enter only through $K$. This is (ii).

\emph{Step 3 (the cost, mass-squared reading).} By clause 2 of the theorem the eigenvalues of $H$ are masses squared. If $\psi$ is as in Step 2 and reproduces the charged-lepton ratios, $\lambda_{e} \ge 1 + x$ and $\lambda_{\tau} = (m_{\tau}/m_{e})^{2}\lambda_{e} \ge (m_{\tau}/m_{e})^{2}(1 + x)$. With $m_{e} = 0.51099895$ MeV, $m_{\mu} = 105.6583755$ MeV and $m_{\tau} \in (1776.84, 1777.02)$ MeV (the store's PDG 2024 values), $(m_{\tau}/m_{e})^{2} \in (1.2090, 1.2094) \times 10^{7}$ and $(m_{\mu}/m_{e})^{2} \in (42753.0, 42753.2)$. At $x^{*} \in (769.3324, 769.3326)$ (Proposition \claimref{families.koide-class-crossing-fit-conclusion}) this gives $\lambda_{\tau} > 9.313 \times 10^{9}$, while $1 + 6x^{*} + 6x^{*2} < 3.556 \times 10^{6}$. Write $A = (1 + x)\mathbb{1} + (5x + 4x^{2})P$, with eigenvalues $1 + 6x + 4x^{2}, 1 + x, 1 + x$. Weyl's inequalities $\lambda_{i + j - 1}(A + K) \le \lambda_{i}(A) + \lambda_{j}(K)$ (eigenvalues in decreasing order) give $\lambda_{1}(K) \ge \lambda_{\tau} - (1 + 6x + 4x^{2})$ and $\lambda_{2}(K) \ge \lambda_{\mu} - (1 + 6x + 4x^{2})$. With $\lambda_{\tau} > 9.313 \times 10^{9}$ the first exceeds $7{,}860 \cdot 2x^{*2}$, and with $\lambda_{\mu} \ge (m_{\mu}/m_{e})^{2}(1 + x^{*})$ the second exceeds $3.05 \times 10^{7}$, where the skeleton's $K = 2x^{2}P$ has rank one (the check brackets all of these exactly). Neither is a correction to the skeleton. This is (iii).

\emph{Scope.} The statement is kinematic and exact. It fixes the second moments at the skeleton's; it does not exclude completions that move the second moments, and by (i) such a completion lowers the electron below $\cosh^{2}r$ only by concentrating the face quanta in one family direction, the electron's. Within the displaced-Gaussian class, free means achieve such a concentration only as the relabelling of Proposition \claimref{families.classification-displaced-gaussian-class-solved-level}(3). The mass-squared reading of $H$ is the skeleton theorem's; under a linear reading the factors in (iii) change but (i) and (ii) do not.
